\documentclass[12pt,a4paper,openany,oneside]{book}

% --- Paquetes ---
\usepackage[utf8]{inputenc}
\usepackage[spanish, es-noshorthands]{babel}
\usepackage{amsmath, amsfonts, amssymb}
\usepackage{graphicx}
\usepackage{geometry}
\usepackage{hyperref}
\usepackage{booktabs}
\usepackage{fancyhdr}
\usepackage{listings}
\usepackage{xcolor}
\usepackage{tikz}
\usepackage{pgfplots}
\usepackage{colortbl}
\usepackage{multirow}
\usepackage{hhline}
\usetikzlibrary{arrows.meta, positioning, shapes.geometric, calc, shapes, chains, scopes, shadows}
\pgfplotsset{compat=1.18}

\geometry{margin=2.5cm}

% --- Colores y estilos para código Python ---
\definecolor{codegreen}{rgb}{0,0.6,0}
\definecolor{codegray}{rgb}{0.5,0.5,0.5}
\definecolor{codepurple}{rgb}{0.58,0,0.82}
\definecolor{backcolour}{rgb}{0.95,0.95,0.92}

\lstset{
    language=Python,
    backgroundcolor=\color{backcolour},   
    commentstyle=\color{codegreen},
    keywordstyle=\color{blue},
    numberstyle=\tiny\color{codegray},
    stringstyle=\color{codepurple},
    basicstyle=\ttfamily\small,
    breaklines=true,
    frame=single,
    showstringspaces=false
}

% --- Estilo de cabecera y pie ---
\pagestyle{fancy}
\fancyhf{}
\renewcommand{\headrulewidth}{0.4pt}
\renewcommand{\footrulewidth}{0.4pt}
\fancyhead[L]{\small Carlos César Sánchez Coronel}
\fancyhead[R]{\small Sesión 14: Despliegue de Modelos NLP (MLOps)}
\fancyfoot[L]{\small Carlos César Sánchez Coronel}
\fancyfoot[C]{\thepage}

% --- Portada ---
\title{
    \textbf{Sesión 14} \\ 
    \Huge \textbf{DESPLIEGUE DE MODELOS NLP} \\
    \Large MLOps, APIs, Contenerización y Monitoreo
}
\author{
    \small Por \\ \vspace{0.2cm}
    \Large \textsc{Carlos César Sánchez Coronel}
}
\date{2026}

\begin{document}

\maketitle
\tableofcontents
\addcontentsline{toc}{chapter}{Índice general}

% ------------------------------------------------------------------
% CAPÍTULO 14: DESPLIEGUE DE MODELOS NLP (MLOPS)
% ------------------------------------------------------------------
\chapter{Despliegue de Modelos NLP (MLOps)}

El despliegue es el proceso de llevar un modelo de NLP desde un entorno de experimentación (como un notebook) a un entorno de producción donde los usuarios reales puedan interactuar con él. En el caso de los sistemas de NLP modernos, esto implica no solo servir el modelo, sino gestionar la infraestructura de un Chatbot completo, incluyendo latencia, escalabilidad y persistencia de datos.

\section{Logro de la sesión}
Llevar un modelo NLP a producción, crear APIs robustas, diseñar la arquitectura de un chatbot (Frontend/Backend) y monitorear su rendimiento en tiempo real.

\section{Arquitectura de un Chatbot de Producción}

Un despliegue profesional separa el entrenamiento de la inferencia. El flujo típico sigue una estructura multicapa:

\begin{center}
\begin{tikzpicture}[
    node distance=1.2cm,
    layer/.style={rectangle, draw, fill=blue!5, text width=3cm, align=center, minimum height=1cm, rounded corners, font=\small, drop shadow},
    database/.style={cylinder, draw, shape border rotate=90, fill=green!10, text width=1.5cm, align=center, minimum height=1.5cm, font=\tiny, drop shadow},
    arrow/.style={-Stealth, thick}
]
    % Nodos
    \node[layer] (front) {Frontend \\ (React / Streamlit)};
    \node[layer, right=of front] (api) {Backend API \\ (FastAPI / Flask)};
    \node[layer, right=of api, fill=orange!10] (model) {Inferencia \\ (ONNX / Triton)};
    \node[database, below=of api] (redis) {Memoria Sesión \\ (Redis)};
    \node[database, below=of model] (vectordb) {Knowledge Base \\ (Chroma/Pinecone)};

    % Conexiones
    \draw[arrow] (front) -- (api);
    \draw[arrow] (api) -- (model);
    \draw[arrow] (api) -- (redis);
    \draw[arrow] (model) -- (vectordb);
    \draw[arrow] (api.west) -- ++(-0.5,0) |- (front.west) node[pos=0.2, left] {Respuesta};

\end{tikzpicture}
\captionof{figure}{Arquitectura desacoplada de un chatbot NLP empresarial.}
\end{center}

\section{Serialización y Formatos de Exportación}

Para que un modelo sea eficiente en producción, debe serializarse. Los formatos comunes incluyen:
\begin{itemize}
    \item \textbf{Pickle / Joblib}: Formato nativo de Python, fácil de usar pero con riesgos de seguridad y poca portabilidad.
    \item \textbf{ONNX (Open Neural Network Exchange)}: Optimiza el grafo del modelo para una ejecución más rápida en diversos hardwares.
    \item \textbf{TorchScript}: Permite ejecutar modelos de PyTorch en entornos de C++ sin dependencia de Python.
\end{itemize}

\section{Backend: Creación de APIs con FastAPI}

\textbf{FastAPI} se ha convertido en el estándar para NLP debido a su manejo asíncrono, lo cual es vital cuando el modelo tarda cientos de milisegundos en generar una respuesta.

\begin{lstlisting}[language=Python, caption=Ejemplo de API básica para Chatbot]
from fastapi import FastAPI
from transformers import pipeline

app = FastAPI()
chatbot = pipeline("text-generation", model="gpt2")

@app.post("/chat")
async def ask_chatbot(prompt: str):
    # Logica de inferencia asincrona
    response = chatbot(prompt, max_length=50)
    return {"reply": response[0]['generated_text']}
\end{lstlisting}

\section{Contenerización con Docker}

El uso de **Docker** garantiza que el modelo se ejecute en el mismo entorno (librerías de CUDA, versiones de transformadores) sin importar el servidor de despliegue.

\begin{itemize}
    \item \textbf{Dockerfile}: Archivo de configuración que define la imagen.
    \item \textbf{Docker Compose}: Permite levantar el Chatbot, la Base de Datos Vectorial y el Backend con un solo comando.
\end{itemize}

\section{Monitoreo y MLOps: Observabilidad}

Una vez el modelo está "vivo", debemos vigilar su salud mediante métricas de rendimiento y de datos:

\begin{table}[h]
\centering
\begin{tabular}{|l|p{8cm}|}
\hline
\rowcolor{gray!20} \textbf{Métrica} & \textbf{Descripción} \\ \hline
\textbf{Latencia} & Tiempo transcurrido desde el prompt hasta la respuesta. \\ \hline
\textbf{Throughput} & Número de solicitudes que el sistema maneja por segundo. \\ \hline
\textbf{Data Drift} & Cambio en las características estadísticas de las preguntas del usuario respecto al entrenamiento. \\ \hline
\textbf{Concept Drift} & Cambio en la relación entre la entrada y la salida (el modelo ya no predice correctamente la intención). \\ \hline
\end{tabular}
\caption{Métricas críticas de monitoreo en producción.}
\end{table}

\section{Preparación para Entrevistas}

\subsection{Pregunta Clave: "¿Cómo desplegarías un modelo NLP en producción?"}
\textbf{Respuesta sugerida}: Utilizaría una separación clara entre entrenamiento e inferencia. Exportaría el modelo a **ONNX** para optimizar la velocidad. Construiría una API con **FastAPI** para manejar peticiones asíncronas y empaquetaría todo en un contenedor **Docker**. Finalmente, implementaría monitoreo de **latencia** y **deriva de datos** (Data Drift) para asegurar la calidad a largo plazo.

\section{Resumen}
\begin{itemize}
    \item La \textbf{serialización} (ONNX, TorchScript) es el primer paso hacia la producción.
    \item El \textbf{Backend} gestiona la lógica, mientras el \textbf{Frontend} ofrece la interfaz de usuario.
    \item El \textbf{versionamiento} (DVC, Hugging Face Hub) es crucial para la reproducibilidad.
    \item El \textbf{monitoreo} constante evita que el modelo quede obsoleto por cambios en el comportamiento del usuario.
\end{itemize}

\end{document}